\documentclass[11pt,a4paper]{article}
\usepackage[T1]{fontenc}
\usepackage{lmodern}
\usepackage[margin=26mm]{geometry}
\usepackage{amsmath,amssymb,amsthm}
\usepackage{microtype}
\usepackage{enumitem}
\usepackage[hidelinks]{hyperref}
\setlength{\parskip}{3pt}
\setlength{\emergencystretch}{3em}
\newtheorem{theorem}{Theorem}[section]
\newtheorem{lemma}[theorem]{Lemma}
\newtheorem{corollary}[theorem]{Corollary}
\theoremstyle{remark}
\newtheorem{remark}[theorem]{Remark}
\newcommand{\Z}{\mathbb Z}
\newcommand{\N}{\mathbb N}
\newcommand{\Q}{\mathbb Q}
\newcommand{\F}{\mathbb F}
\newcommand{\ps}[1]{#1[[x]]}
\hypersetup{pdftitle={A389474: Cubic reversion modulo nine and three},pdfauthor={Dev Goyal}}
\title{A389474\\[3pt]Cubic reversion modulo nine and three}
\author{Dev Goyal}
\date{24 October 2026}
\begin{document}
\maketitle
\begin{center}
\fbox{\begin{minipage}{0.91\linewidth}\small
\textbf{AI contribution and review status.}
ChatGPT (OpenAI) generated the mathematical exposition, proof development,
Lean interfaces, and checks for this collection in response to Dev Goyal's
prompts. Dev Goyal is the sole human author. Independent human review is
pending. Original sequence authors, known mathematical results, and adapted
upstream code retain their separate credits. This document makes no claim
of publication, OEIS acceptance, or worldwide priority.
\end{minipage}}
\end{center}
\begin{abstract}
We construct the normalized integer solution of the defining equation, prove that its square is $x^2+2x^3$ modulo nine, and formalize the associated cubic-power assertion modulo three. The second assertion is a direct application of the classical Frobenius identity.
\end{abstract}
\section{Source and conclusions}
Let $U=x+O(x^2)\in\Z[[x]]$ be specified by
\[
U\bigl(x^2+2x^3+2U(x)^2\bigr)=3U(x)^2.
\]
This is the A389474 source equation \cite{oeis}, with offset one.
We first construct its unique integral solution. We then prove
$U^2\equiv x^2+2x^3\pmod9$ and, for every $k\ge0$,
$U^{3k}\equiv U(x^3)^k\pmod3$. Powers without the composition symbol
mean ordinary multiplication.
\section{Integral construction and uniqueness}

\begin{theorem}\label{thm:cubicexist}
For every $s\in\Z$, there is a unique normalized series
$R_s\in\ps{\Z}$ satisfying
\begin{equation}\label{eq:cubicfamily}
R_s(3x^2)=2x^2+R_s(x)^2+2sR_s(x)^3.
\end{equation}
\end{theorem}

\begin{proof}
For any normalized integral series $U$, define its residual by
\[
E_s(U)=U^2+2sU^3-U(3x^2)+2x^2.
\]
Every coefficient of $E_s(U)$ is even. Indeed, after reduction modulo
two, Frobenius gives $U(x)^2=U(x^2)$ and $3x^2=x^2$, so the residual
vanishes. Also $E_s(U)=O(x^3)$ for every normalized $U$.

Suppose $V-U=cx^r+O(x^{r+1})$, with $r\ge2$. The quadratic difference
has leading term $2cx^{r+1}$; the cubic difference has order at least
$r+2$; and $(V-U)(3x^2)$ has order at least $2r\ge r+2$.
Consequently
\begin{equation}\label{eq:triangular}
E_s(V)-E_s(U)=2cx^{r+1}+O(x^{r+2}).
\end{equation}
Starting with $U=x$, choose its degree-$r$ coefficient to cancel the
degree-$(r+1)$ coefficient of the residual. The needed division by two
is exact by the preceding parity argument. Equation~\eqref{eq:triangular}
shows that earlier cancellations remain valid. Each degree is thus
determined by finitely many earlier coefficients, and the resulting
integral series has zero residual in every degree.

For uniqueness, take the first degree $r\ge2$ where two candidate
solutions differ. Equation~\eqref{eq:triangular} would give a nonzero
degree-$(r+1)$ difference of residuals, because $2c\ne0$ in $\Z$.
This contradicts both residuals being zero.
\end{proof}

The proof formalizes divisibility of the residual by two before using
integer division. Thus the construction does not assume the existence
or integrality of an OEIS sequence. The names used below denote the
uniquely constructed series.

\section{The source bridge and the modulus-nine proof}

Let $R=R_1$ and let $U$ be its compositional inverse. Such an inverse
exists over $\Z$ for every normalized series: its degree-$n$ coefficient
is determined with leading multiplier one, recursively. Both inverse
identities are proved in the reused formal-series infrastructure.
Substituting $U(x)$ in \eqref{eq:cubicfamily} gives
\[
R(3U(x)^2)=2U(x)^2+x^2+2x^3.
\]
Applying $U$ proves the original A389474 equation
\begin{equation}\label{eq:474}
U\bigl(x^2+2x^3+2U(x)^2\bigr)=3U(x)^2.
\end{equation}
Conversely, if a normalized $U$ satisfies \eqref{eq:474}, substitute its
inverse $R$ and then apply $R$ to recover \eqref{eq:cubicfamily} for
$s=1$. Therefore Theorem~\ref{thm:cubicexist} also gives existence and
uniqueness for the literal A389474 source equation.

\begin{theorem}\label{thm:modnine}
The A389474 series satisfies
\[
U(x)^2\equiv x^2+2x^3\pmod9.
\]
In particular, $[x^n]U(x)^2$ is divisible by nine for every $n>3$.
\end{theorem}

\begin{proof}
Work in $(\Z/9\Z)[[x]]$ and put $S=U^2$, $Q=x^2+2x^3$.
The source equation becomes $U(Q+2S)=3S$. Let $V$ be the genuine
normalized compositional inverse of $U$ over this ring. Apply $V$ to
obtain
\[
Q+2S=V(3S).
\]
The series $3S$ has zero constant coefficient and is square-zero:
$(3S)^2=9S^2=0$. A normalized series has the form
$V(t)=t+t^2H(t)$, so $V(3S)=3S$. It follows that $Q+2S=3S$, hence
$S=Q$, as required. No division by three is used.
\end{proof}

This establishes a full formal-series identity, including the low-degree
terms, and hence strengthens the source's statement restricted to
degrees above three.

\section{The standard Frobenius identity}
For any integer formal series $W$ and prime $p$, coefficient reduction gives
\[
\overline W(x)^p=\overline W(x^p)\qquad\text{in }\F_p[[x]].
\]
For each fixed coefficient one may truncate $W$ at that degree. The
polynomial identity $(u+v)^p=u^p+v^p$ in characteristic $p$, together
with $a^p=a$ in $\F_p$, proves the displayed identity coefficient by
coefficient. Thus it holds for every formal series, without a convergence
assumption. This is standard algebra; its application below is not a new
general principle.

Apply the identity at $p=3$ to $W=U$ and raise both sides to the
ordinary $k$th power. For every $n,k\ge0$ this yields
\[
[x^n]\bigl(U(x)^{3k}-U(x^3)^k\bigr)\equiv0\pmod3.
\]
This proves the second selected source assertion, including the boundary
cases $n=0$ and $k=0$. It requires no special equation for $U$.
\section{Lean proof and exact scope}
The dedicated entry file is \path{lean/Entries/A389474.lean}. Its declarations lie in namespace \texttt{FinalCollection.A389474}. The principal checked endpoints are:
\begin{itemize}[nosep]
\item \texttt{source\_equation}: the exact source equation.
\item \texttt{source\_existence\_uniqueness}: existence and uniqueness over the integers.
\item \texttt{integral\_family\_existence}: the entire integer-parameter auxiliary family.
\item \texttt{square\_mod\_nine}: the full modulus-nine identity.
\item \texttt{square\_coefficients}: the source range $n>3$.
\item \texttt{power\_coefficients}: the standard Frobenius consequence for all $n,k\ge0$.
\end{itemize}
The full implementations are in \path{lean/More/CubicConsequences.lean}.
The full integral construction, exact division by two, genuine compositional inversion, and both source-uniqueness directions are present in the shared implementation. No division by three is used modulo nine. Normalization and source uniqueness are explicit endpoints in the separate entry file.

The companion is a separate, explicitly typed Lean entry file. Its proofs
invoke the preserved implementation modules; the complete dependency source
is included in the same archive. Shared lemmas are reused rather than
assumed. The entry file is not a self-contained replacement for Mathlib or
the supplied common library.

The project pins Lean 4.19.0 and Mathlib revision
\begin{center}\small\ttfamily
c44e0c8ee63ca166450922a373c7409c5d26b00b.
\end{center}
The collection checker recompiles the proof sources and audits the
theorems' transitive axiom dependencies. Only \texttt{propext},
\texttt{Classical.choice}, and \texttt{Quot.sound} are permitted.
The delivered verification reports record the actual successful runs.
Finite numerical checks are corroboration; the displayed conclusions are
unbounded theorems. An AI review of source correspondence does not replace
independent human review.
\begin{thebibliography}{9}
\raggedright
\bibitem{oeis} OEIS Foundation, entry A389474,
\url{https://oeis.org/A389474}. Source checked 24 October 2026.
The source sequence and selected conjectures are credited to Paul D. Hanna.
\bibitem{repo} D. Goyal, supplied reference repository,
\url{https://github.com/dvgyl/oeis-a396843-proof}.
This separate presentation belongs to the final collected proof package;
no earlier DOI is reassigned to it.
\end{thebibliography}
\end{document}
